\section{Model results}
\label{sec:results}

Everything below is model output. The comparisons against disclosure come in
Section~\ref{sec:comparison} and the assumption sensitivity in
Section~\ref{sec:sensitivity}. The figures for Schwab have no disclosed counterpart
in either year, which is worth holding in mind throughout.

\subsection{Economic value under the standardised scenarios}

Table~\ref{tab:scenarios} reports the six Basel scenarios for both banks at both
year-ends, on the conservative deposit assumption of a two-year uninsured half-life
and the 0.60 beta SVB discloses.

\begin{table}[htbp]
\centering
\caption{Change in economic value of equity under the six standardised scenarios,
conservative deposit case, per cent of base economic value. Base levels are in \$
millions. All figures are model output; neither bank disclosed an
economic-value profile for 2022 and Schwab disclosed none in either year.}
\label{tab:scenarios}
\small
\begin{tabular}{llrrrrrrr}
\toprule
Bank & Year & Base EVE & Par.\ up & Par.\ down & Short up & Short down & Steep. & Flat. \\
\midrule
SVB    & 2021 & 16,973 & $-62.47$  & $37.73$  & $3.36$    & $-0.38$  & $-48.02$ & $24.43$ \\
SVB    & 2022 & 9,523  & $-93.24$  & $128.75$ & $-0.20$   & $0.22$   & $-66.83$ & $52.87$ \\
Schwab & 2021 & 28,702 & $-89.65$  & $-0.90$  & $-24.24$  & $16.44$  & $-37.82$ & $-6.62$ \\
Schwab & 2022 & 10,862 & $-161.83$ & $172.01$ & $-52.26$  & $53.96$  & $-56.08$ & $20.92$ \\
\bottomrule
\end{tabular}
\end{table}

Two cautions before reading it.

The percentages are taken against base economic value, and base economic value falls
sharply over 2022 for both banks: SVB from \$16,973m to \$9,523m, Schwab from
\$28,702m to \$10,862m. A given dollar move therefore registers as a much larger
percentage in the 2022 rows than in the 2021 rows, and the $-161.83\%$ in Schwab's
corner is a small denominator as much as a large numerator. Comparisons down a column
across years are not like-for-like; comparisons across a row are.

The conservative deposit case is also a harsher assumption for Schwab than for SVB. A
two-year uninsured half-life applied to a base that was 18.8\% uninsured is a
stringent reading of a funding structure that was mostly insured and mostly sticky,
whereas the same assumption applied to SVB's 96\% uninsured base is closer to a
plausible central case. The conservative column is included for both banks so that
one framework is applied to both, not because the assumption fits both equally. This
is why the comparison that carries the paper's argument is the one in
Section~\ref{sec:sensitivity}, where the deposit assumption is varied rather than
fixed.

With those in place, the pattern is as expected in direction and larger than expected
in size. Both banks lose heavily in a parallel rise and in a steepener, which are the
two scenarios that mark down a long asset book. Both are helped by a parallel fall at
the end of 2022, when there was room beneath the curve for the shock to work. Neither
is much moved by the short-rate scenarios at the end of 2021, when the front of the
curve was already at the floor and a downward short shock had almost nothing to
absorb. SVB's $-0.38\%$ at short down is the floor binding, not an absence of
exposure.

Schwab at the end of 2021 is the one row that does not fit the pattern. It loses value
in a parallel fall ($-0.90\%$) and in a flattener ($-6.62\%$), where SVB gains in
both. On this model that is the deposit franchise dominating: a bank funded largely by
insured deposits paying close to nothing earns a wide spread on those balances, and a
fall in rates compresses it faster than the asset repricing compensates. It is worth
flagging as a model-dependent result rather than a finding about Schwab, because it is
precisely the quantity that no disclosure constrains.

\subsection{The supervisory outlier test}

The Basel framework treats a bank whose worst-case economic-value loss exceeds 15\% of
Tier~1 capital as warranting supervisory attention \citep{BCBS2016}.
Table~\ref{tab:outlier} applies that screen.

\begin{table}[htbp]
\centering
\caption{Supervisory outlier test. Worst of the six standardised scenarios, measured
against consolidated Tier~1 capital as reported for each date. The 2021 figures are
taken from the comparative column of each bank's FY2022 filing, so that both years
rest on one source; SVB's FY2021 filing states its own Tier~1 capital \$1m higher,
which moves the ratio by 0.002 percentage points. The implied-by-disclosure row uses
the deposit half-life solved from SVB's own published sensitivity; the conservative
rows use the two-year half-life.}
\label{tab:outlier}
\small
\begin{tabular}{llrrrc}
\toprule
Bank & Year & Deposit case & Worst loss (\%) & Loss / Tier 1 (\%) & Breaches 15\% \\
\midrule
SVB    & 2021 & implied by disclosure & $-27.61$  & 36.24 & yes \\
SVB    & 2021 & conservative          & $-62.47$  & 65.43 & yes \\
SVB    & 2022 & conservative          & $-93.24$  & 50.73 & yes \\
Schwab & 2021 & conservative          & $-89.65$  & 67.86 & yes \\
Schwab & 2022 & conservative          & $-161.83$ & 43.62 & yes \\
\bottomrule
\end{tabular}
\end{table}

Every bank-year breaches the threshold, and the breach is not marginal in any of them:
the smallest ratio is 36.24\% against a 15\% screen.

The first row is the one that matters most, because it is the only row with external
support. It uses the deposit half-life implied by SVB's own disclosure, so the
exposure it reports is the one consistent with what SVB published. On that basis
SVB's worst-case loss at the end of 2021 was 36.24\% of Tier~1 capital, more than
twice the supervisory threshold, and that figure is derived entirely from
information in the FY2021 filing, well before the failure. Nothing in
it required hindsight, a view on the path of rates, or any data the filing did not
contain.

That every bank-year breaches, Schwab's included, is the reason the test cannot carry
the comparison on its own. The screen separates banks with large interest-rate
exposure from banks without it, and in 2021 and 2022 both of these banks had large
exposure. What it does not capture is the funding structure that decides whether such
exposure becomes a failure, which is the subject of the next subsection.

\subsection{Run thresholds}

The run model asks a different question: what fraction of uninsured deposits has to
leave before the bank cannot meet the demand. The question is posed in the tradition
of \citet{DiamondDybvig1983}, in which a bank funding illiquid assets with
on-demand claims admits an equilibrium where withdrawal is individually rational
because others are withdrawing; what is computed here is not that equilibrium but the
balance-sheet threshold beyond which the demand cannot be met at all. Outflows are
met from cash first, then
secured borrowing, then sales of available-for-sale securities, then sales of
held-to-maturity securities, with the accounting consequence that selling from the
held-to-maturity portfolio forces the remainder to be marked.

Two distinct failure points come out of that waterfall. A solvency threshold, where
realised losses exhaust equity, and a liquidity threshold, where the bank runs out of
assets it can turn into cash regardless of its accounting equity. The binding
threshold is whichever comes first.

For SVB at the end of 2022 the binding threshold is \textbf{liquidity, at 67.5\%} of
uninsured deposits. The solvency threshold is never reached in any of the four
bank-years: the constraint that binds is the inability to raise cash, not the
exhaustion of book equity. At a full run on all uninsured balances, \$54,043m of
demand would have gone unfunded, and the fire-sale discount at which the solvency
constraint would have become the binding one is 28.3\%.

The 2021 threshold is 73.0\%, so the model has the bank becoming more fragile over
2022, driven by the held-to-maturity loss rising from \$968m to \$15,152m, which
takes the after-tax consequence of recognising it from 4.7\% to 74.8\% of book
equity.

Set against that, SVB lost over \$40bn of deposits on 9~March 2023 and management
expected to lose over \$100bn more the following day \citep{FRB2023}. Taken together
that is \$140bn of demand across the two days, or 84.6\% of the uninsured base
measured at the end of 2022, against a modelled threshold of 67.5\%.

The agreement in order of magnitude is worth noting and easy to over-read, so three
qualifications. The model is a static balance-sheet calculation with no account of the
speed at which a modern run moves, which makes 67.5\% better understood as an upper
bound on what the balance sheet could absorb than as a forecast of what would trigger
failure. The \$100bn is an expectation recorded by management, not a completed
withdrawal: the bank was closed on 10~March, so the two-day figure measures demand
rather than settled outflow. And the denominator is the uninsured balance at
31~December 2022, which had been declining; the base on the day was not the base in
the model. None of these is a reason to discard the comparison, but it is a comparison
between a modelled bound and a reported demand, not between a prediction and an
outcome.

Schwab crosses neither threshold in either year. That holds across all 42
combinations of secured borrowing access and fire-sale discount tried in
Section~\ref{sec:sensitivity}: in every one, cash, borrowing capacity and
available-for-sale securities together cover a complete run on uninsured balances,
and nothing is left unfunded. The two banks carried comparable unrecognised losses
into the end of 2022, \$15,152m against \$14,138m, and on this model only one of
them had a funding structure that could turn those losses into a failure.
